% !PRINTED 627
\weberpage{627}

%<<CONT
and if one here puts rational quantities for $\lambda$ and $\mu$ (from
an arbitrary domain of rationality, e.g. also from the field
of rational functions of $\lambda$ and $\mu$), then the equation
of degree 5
\weq{24}{x^5+\alpha x+\beta=0}
is algebraically solvable, since indeed the case that this equation
becomes reducible also leads to algebraically solvable equations.

Conversely we can say that no irreducible equation
of degree 5 of the form (24) is algebraically solvable in
which the coefficients $\alpha$, $\beta$ are not representable in the form (23).

If we set, e.g., $\lambda=-1$, $\mu=1$, we obtain $64\alpha=-5^4$,
$64\beta=-5^4$. The equation (24) is simplified, if we introduce for
$x$ a new unknown $\xi$, by the equation $x\xi=5$.
We thus obtain the equation
\[
\xi^5+5\xi^4-5\cdot 64=0
\]
as an example of an \gesperrt{algebraically solvable} equation which, moreover,
is \gesperrt{irreducible} by the theorem mentioned above.

\begin{center}
\S.\,182.

\gesperrt{The Group of the Resolvent.}
\end{center}

The theorems that we have derived in the preceding section
admit a remarkable conclusion about the equations
of degree 6.

We now again take the $x_0$, $x_1$, $x_2$, $x_3$, $x_4$ of the preceding
section as independent variables. Then the six
quantities
\weq{1}{\begin{aligned}
v_1&=(u_1-u_1')^2, & v_2&=(u_2-u_2')^2, & v_3&=(u_3-u_3')^2\\
v_4&=(u_4-u_4')^2, & v_5&=(u_5-u_5')^2, & v_6&=(u_6-u_6')^2
\end{aligned}}
are the roots of an equation of degree 6, $F(v)=0$, whose
coefficients are rationally expressible in terms of the symmetric fundamental functions
$a$, $b$, $c$, $d$, $e$ of the $x$, and which is irreducible in the field
of the rational functions of these quantities. This
equation is a resolvent of the equation of degree 5, $f(x)=0$,
whose roots are the $x$.

If $M$ is the full metacyclic group of the $x$, then the
groups $\pi^{-1}M\pi$ conjugate to $M$ have no other common
%CONT>>
